% Môn: Toán
% Lớp: 10
% Chương: Chương I. Mệnh đề và tập hợp
% Bài: Bài 2. Tập hợp và các phép toán trên tập hợp
% Loại: Lý thuyết
% File riêng, không trộn vào de.tex
% Định nghĩa lệnh Lý thuyết — dùng khi biên dịch lt.tex bằng TeX.
% App web đọc lt.tex trực tiếp (bỏ \input file này).
% Trong lt.tex, từ thư mục bài:
%   % Định nghĩa lệnh Lý thuyết — dùng khi biên dịch lt.tex bằng TeX.
% App web đọc lt.tex trực tiếp (bỏ \input file này).
% Trong lt.tex, từ thư mục bài:
%   % Định nghĩa lệnh Lý thuyết — dùng khi biên dịch lt.tex bằng TeX.
% App web đọc lt.tex trực tiếp (bỏ \input file này).
% Trong lt.tex, từ thư mục bài:
%   \input{../../../../_lenh/lythuyet.tex}
% từ gốc repo:
%   \input{ngan-hang/_lenh/lythuyet.tex}
\makeatletter
\providecommand{\indam}[1]{\textbf{#1}}
\providecommand{\chuy}[1]{\par\smallskip\noindent\textit{#1}\par}
\providecommand{\luuy}[1]{\par\smallskip\noindent\textbf{Lưu ý.} #1\par}
\providecommand{\ghichu}[1]{\par\smallskip\noindent\textbf{Ghi chú.} #1\par}
\providecommand{\vidu}[1]{\par\smallskip\noindent\textbf{Ví dụ.} #1\par}
\providecommand{\Blythuyet}{\subsection*{Lý thuyết}}
% Hai cột: kiến thức | hình
\providecommand{\haicotchay}[2]{%
  \par\medskip\noindent
  \begin{minipage}[t]{0.52\linewidth}#1\end{minipage}\hfill
  \begin{minipage}[t]{0.46\linewidth}#2\end{minipage}%
  \par\medskip
}
\providecommand{\kienthuccot}[2]{\haicotchay{#1}{#2}}
% Dự phòng nếu chưa nạp graphicx / caption (không ghi đè bản thật)
\providecommand{\resizebox}[3]{#3}
\providecommand{\captionof}[2]{\par\smallskip\noindent\textit{#2}\par}
\newenvironment{hoatdong}[1]{%
  \par\medskip\noindent\textbf{Hoạt động.} #1\par\smallskip
}{\par\medskip}
\newenvironment{emcobiet}{%
  \par\medskip\noindent\textbf{Em có biết}\par\smallskip
}{\par\medskip}
\newenvironment{emdahoc}{%
  \par\medskip\noindent\textbf{Em đã học}\par\smallskip
}{\par\medskip}
\newenvironment{emcothe}{%
  \par\medskip\noindent\textbf{Em có thể}\par\smallskip
}{\par\medskip}
\newenvironment{kienthuc}{%
  \par\medskip\noindent\textbf{Kiến thức cốt lõi}\par\smallskip
}{\par\medskip}
\newenvironment{traloi}{%
  \par\smallskip\noindent\textit{Trả lời / ghi chú của em}\par
  \vspace{1.2em}%
}{\par\medskip}
\makeatother

% từ gốc repo:
%   % Định nghĩa lệnh Lý thuyết — dùng khi biên dịch lt.tex bằng TeX.
% App web đọc lt.tex trực tiếp (bỏ \input file này).
% Trong lt.tex, từ thư mục bài:
%   \input{../../../../_lenh/lythuyet.tex}
% từ gốc repo:
%   \input{ngan-hang/_lenh/lythuyet.tex}
\makeatletter
\providecommand{\indam}[1]{\textbf{#1}}
\providecommand{\chuy}[1]{\par\smallskip\noindent\textit{#1}\par}
\providecommand{\luuy}[1]{\par\smallskip\noindent\textbf{Lưu ý.} #1\par}
\providecommand{\ghichu}[1]{\par\smallskip\noindent\textbf{Ghi chú.} #1\par}
\providecommand{\vidu}[1]{\par\smallskip\noindent\textbf{Ví dụ.} #1\par}
\providecommand{\Blythuyet}{\subsection*{Lý thuyết}}
% Hai cột: kiến thức | hình
\providecommand{\haicotchay}[2]{%
  \par\medskip\noindent
  \begin{minipage}[t]{0.52\linewidth}#1\end{minipage}\hfill
  \begin{minipage}[t]{0.46\linewidth}#2\end{minipage}%
  \par\medskip
}
\providecommand{\kienthuccot}[2]{\haicotchay{#1}{#2}}
% Dự phòng nếu chưa nạp graphicx / caption (không ghi đè bản thật)
\providecommand{\resizebox}[3]{#3}
\providecommand{\captionof}[2]{\par\smallskip\noindent\textit{#2}\par}
\newenvironment{hoatdong}[1]{%
  \par\medskip\noindent\textbf{Hoạt động.} #1\par\smallskip
}{\par\medskip}
\newenvironment{emcobiet}{%
  \par\medskip\noindent\textbf{Em có biết}\par\smallskip
}{\par\medskip}
\newenvironment{emdahoc}{%
  \par\medskip\noindent\textbf{Em đã học}\par\smallskip
}{\par\medskip}
\newenvironment{emcothe}{%
  \par\medskip\noindent\textbf{Em có thể}\par\smallskip
}{\par\medskip}
\newenvironment{kienthuc}{%
  \par\medskip\noindent\textbf{Kiến thức cốt lõi}\par\smallskip
}{\par\medskip}
\newenvironment{traloi}{%
  \par\smallskip\noindent\textit{Trả lời / ghi chú của em}\par
  \vspace{1.2em}%
}{\par\medskip}
\makeatother

\makeatletter
\providecommand{\indam}[1]{\textbf{#1}}
\providecommand{\chuy}[1]{\par\smallskip\noindent\textit{#1}\par}
\providecommand{\luuy}[1]{\par\smallskip\noindent\textbf{Lưu ý.} #1\par}
\providecommand{\ghichu}[1]{\par\smallskip\noindent\textbf{Ghi chú.} #1\par}
\providecommand{\vidu}[1]{\par\smallskip\noindent\textbf{Ví dụ.} #1\par}
\providecommand{\Blythuyet}{\subsection*{Lý thuyết}}
% Hai cột: kiến thức | hình
\providecommand{\haicotchay}[2]{%
  \par\medskip\noindent
  \begin{minipage}[t]{0.52\linewidth}#1\end{minipage}\hfill
  \begin{minipage}[t]{0.46\linewidth}#2\end{minipage}%
  \par\medskip
}
\providecommand{\kienthuccot}[2]{\haicotchay{#1}{#2}}
% Dự phòng nếu chưa nạp graphicx / caption (không ghi đè bản thật)
\providecommand{\resizebox}[3]{#3}
\providecommand{\captionof}[2]{\par\smallskip\noindent\textit{#2}\par}
\newenvironment{hoatdong}[1]{%
  \par\medskip\noindent\textbf{Hoạt động.} #1\par\smallskip
}{\par\medskip}
\newenvironment{emcobiet}{%
  \par\medskip\noindent\textbf{Em có biết}\par\smallskip
}{\par\medskip}
\newenvironment{emdahoc}{%
  \par\medskip\noindent\textbf{Em đã học}\par\smallskip
}{\par\medskip}
\newenvironment{emcothe}{%
  \par\medskip\noindent\textbf{Em có thể}\par\smallskip
}{\par\medskip}
\newenvironment{kienthuc}{%
  \par\medskip\noindent\textbf{Kiến thức cốt lõi}\par\smallskip
}{\par\medskip}
\newenvironment{traloi}{%
  \par\smallskip\noindent\textit{Trả lời / ghi chú của em}\par
  \vspace{1.2em}%
}{\par\medskip}
\makeatother

% từ gốc repo:
%   % Định nghĩa lệnh Lý thuyết — dùng khi biên dịch lt.tex bằng TeX.
% App web đọc lt.tex trực tiếp (bỏ \input file này).
% Trong lt.tex, từ thư mục bài:
%   % Định nghĩa lệnh Lý thuyết — dùng khi biên dịch lt.tex bằng TeX.
% App web đọc lt.tex trực tiếp (bỏ \input file này).
% Trong lt.tex, từ thư mục bài:
%   \input{../../../../_lenh/lythuyet.tex}
% từ gốc repo:
%   \input{ngan-hang/_lenh/lythuyet.tex}
\makeatletter
\providecommand{\indam}[1]{\textbf{#1}}
\providecommand{\chuy}[1]{\par\smallskip\noindent\textit{#1}\par}
\providecommand{\luuy}[1]{\par\smallskip\noindent\textbf{Lưu ý.} #1\par}
\providecommand{\ghichu}[1]{\par\smallskip\noindent\textbf{Ghi chú.} #1\par}
\providecommand{\vidu}[1]{\par\smallskip\noindent\textbf{Ví dụ.} #1\par}
\providecommand{\Blythuyet}{\subsection*{Lý thuyết}}
% Hai cột: kiến thức | hình
\providecommand{\haicotchay}[2]{%
  \par\medskip\noindent
  \begin{minipage}[t]{0.52\linewidth}#1\end{minipage}\hfill
  \begin{minipage}[t]{0.46\linewidth}#2\end{minipage}%
  \par\medskip
}
\providecommand{\kienthuccot}[2]{\haicotchay{#1}{#2}}
% Dự phòng nếu chưa nạp graphicx / caption (không ghi đè bản thật)
\providecommand{\resizebox}[3]{#3}
\providecommand{\captionof}[2]{\par\smallskip\noindent\textit{#2}\par}
\newenvironment{hoatdong}[1]{%
  \par\medskip\noindent\textbf{Hoạt động.} #1\par\smallskip
}{\par\medskip}
\newenvironment{emcobiet}{%
  \par\medskip\noindent\textbf{Em có biết}\par\smallskip
}{\par\medskip}
\newenvironment{emdahoc}{%
  \par\medskip\noindent\textbf{Em đã học}\par\smallskip
}{\par\medskip}
\newenvironment{emcothe}{%
  \par\medskip\noindent\textbf{Em có thể}\par\smallskip
}{\par\medskip}
\newenvironment{kienthuc}{%
  \par\medskip\noindent\textbf{Kiến thức cốt lõi}\par\smallskip
}{\par\medskip}
\newenvironment{traloi}{%
  \par\smallskip\noindent\textit{Trả lời / ghi chú của em}\par
  \vspace{1.2em}%
}{\par\medskip}
\makeatother

% từ gốc repo:
%   % Định nghĩa lệnh Lý thuyết — dùng khi biên dịch lt.tex bằng TeX.
% App web đọc lt.tex trực tiếp (bỏ \input file này).
% Trong lt.tex, từ thư mục bài:
%   \input{../../../../_lenh/lythuyet.tex}
% từ gốc repo:
%   \input{ngan-hang/_lenh/lythuyet.tex}
\makeatletter
\providecommand{\indam}[1]{\textbf{#1}}
\providecommand{\chuy}[1]{\par\smallskip\noindent\textit{#1}\par}
\providecommand{\luuy}[1]{\par\smallskip\noindent\textbf{Lưu ý.} #1\par}
\providecommand{\ghichu}[1]{\par\smallskip\noindent\textbf{Ghi chú.} #1\par}
\providecommand{\vidu}[1]{\par\smallskip\noindent\textbf{Ví dụ.} #1\par}
\providecommand{\Blythuyet}{\subsection*{Lý thuyết}}
% Hai cột: kiến thức | hình
\providecommand{\haicotchay}[2]{%
  \par\medskip\noindent
  \begin{minipage}[t]{0.52\linewidth}#1\end{minipage}\hfill
  \begin{minipage}[t]{0.46\linewidth}#2\end{minipage}%
  \par\medskip
}
\providecommand{\kienthuccot}[2]{\haicotchay{#1}{#2}}
% Dự phòng nếu chưa nạp graphicx / caption (không ghi đè bản thật)
\providecommand{\resizebox}[3]{#3}
\providecommand{\captionof}[2]{\par\smallskip\noindent\textit{#2}\par}
\newenvironment{hoatdong}[1]{%
  \par\medskip\noindent\textbf{Hoạt động.} #1\par\smallskip
}{\par\medskip}
\newenvironment{emcobiet}{%
  \par\medskip\noindent\textbf{Em có biết}\par\smallskip
}{\par\medskip}
\newenvironment{emdahoc}{%
  \par\medskip\noindent\textbf{Em đã học}\par\smallskip
}{\par\medskip}
\newenvironment{emcothe}{%
  \par\medskip\noindent\textbf{Em có thể}\par\smallskip
}{\par\medskip}
\newenvironment{kienthuc}{%
  \par\medskip\noindent\textbf{Kiến thức cốt lõi}\par\smallskip
}{\par\medskip}
\newenvironment{traloi}{%
  \par\smallskip\noindent\textit{Trả lời / ghi chú của em}\par
  \vspace{1.2em}%
}{\par\medskip}
\makeatother

\makeatletter
\providecommand{\indam}[1]{\textbf{#1}}
\providecommand{\chuy}[1]{\par\smallskip\noindent\textit{#1}\par}
\providecommand{\luuy}[1]{\par\smallskip\noindent\textbf{Lưu ý.} #1\par}
\providecommand{\ghichu}[1]{\par\smallskip\noindent\textbf{Ghi chú.} #1\par}
\providecommand{\vidu}[1]{\par\smallskip\noindent\textbf{Ví dụ.} #1\par}
\providecommand{\Blythuyet}{\subsection*{Lý thuyết}}
% Hai cột: kiến thức | hình
\providecommand{\haicotchay}[2]{%
  \par\medskip\noindent
  \begin{minipage}[t]{0.52\linewidth}#1\end{minipage}\hfill
  \begin{minipage}[t]{0.46\linewidth}#2\end{minipage}%
  \par\medskip
}
\providecommand{\kienthuccot}[2]{\haicotchay{#1}{#2}}
% Dự phòng nếu chưa nạp graphicx / caption (không ghi đè bản thật)
\providecommand{\resizebox}[3]{#3}
\providecommand{\captionof}[2]{\par\smallskip\noindent\textit{#2}\par}
\newenvironment{hoatdong}[1]{%
  \par\medskip\noindent\textbf{Hoạt động.} #1\par\smallskip
}{\par\medskip}
\newenvironment{emcobiet}{%
  \par\medskip\noindent\textbf{Em có biết}\par\smallskip
}{\par\medskip}
\newenvironment{emdahoc}{%
  \par\medskip\noindent\textbf{Em đã học}\par\smallskip
}{\par\medskip}
\newenvironment{emcothe}{%
  \par\medskip\noindent\textbf{Em có thể}\par\smallskip
}{\par\medskip}
\newenvironment{kienthuc}{%
  \par\medskip\noindent\textbf{Kiến thức cốt lõi}\par\smallskip
}{\par\medskip}
\newenvironment{traloi}{%
  \par\smallskip\noindent\textit{Trả lời / ghi chú của em}\par
  \vspace{1.2em}%
}{\par\medskip}
\makeatother

\makeatletter
\providecommand{\indam}[1]{\textbf{#1}}
\providecommand{\chuy}[1]{\par\smallskip\noindent\textit{#1}\par}
\providecommand{\luuy}[1]{\par\smallskip\noindent\textbf{Lưu ý.} #1\par}
\providecommand{\ghichu}[1]{\par\smallskip\noindent\textbf{Ghi chú.} #1\par}
\providecommand{\vidu}[1]{\par\smallskip\noindent\textbf{Ví dụ.} #1\par}
\providecommand{\Blythuyet}{\subsection*{Lý thuyết}}
% Hai cột: kiến thức | hình
\providecommand{\haicotchay}[2]{%
  \par\medskip\noindent
  \begin{minipage}[t]{0.52\linewidth}#1\end{minipage}\hfill
  \begin{minipage}[t]{0.46\linewidth}#2\end{minipage}%
  \par\medskip
}
\providecommand{\kienthuccot}[2]{\haicotchay{#1}{#2}}
% Dự phòng nếu chưa nạp graphicx / caption (không ghi đè bản thật)
\providecommand{\resizebox}[3]{#3}
\providecommand{\captionof}[2]{\par\smallskip\noindent\textit{#2}\par}
\newenvironment{hoatdong}[1]{%
  \par\medskip\noindent\textbf{Hoạt động.} #1\par\smallskip
}{\par\medskip}
\newenvironment{emcobiet}{%
  \par\medskip\noindent\textbf{Em có biết}\par\smallskip
}{\par\medskip}
\newenvironment{emdahoc}{%
  \par\medskip\noindent\textbf{Em đã học}\par\smallskip
}{\par\medskip}
\newenvironment{emcothe}{%
  \par\medskip\noindent\textbf{Em có thể}\par\smallskip
}{\par\medskip}
\newenvironment{kienthuc}{%
  \par\medskip\noindent\textbf{Kiến thức cốt lõi}\par\smallskip
}{\par\medskip}
\newenvironment{traloi}{%
  \par\smallskip\noindent\textit{Trả lời / ghi chú của em}\par
  \vspace{1.2em}%
}{\par\medskip}
\makeatother


\subsection{Lý thuyết}

\subsubsection{Các khái niệm cơ bản về tập hợp}
\paragaraph{Tập Hợp}
\begin{hoatdong}{Hoạt động 1}
Trong một buổi sinh hoạt chuyên đề, có hai chuyên đề được tổ chức. Chuyên đề 1 có các thành viên: Nam, Hương, Chi, Tú, Bình, Ngân, Khánh. Chuyên đề 2 có các thành viên: Hương, Chi, Tú, Khánh, Bình, Hân, Hiền, Lam.
Gọi $A$ là tập hợp những thành viên tham gia Chuyên đề 1, $B$ là tập hợp những thành viên tham gia Chuyên đề 2.
\begin{itemize}
    \item[a)] Nam có là một phần tử của tập hợp $A$ không? Ngân có là một phần tử của tập hợp $B$ không?
    \item[b)] Hãy mô tả các tập hợp $A$ và $B$ bằng cách liệt kê các phần tử.
\end{itemize}
\end{hoatdong}
\loigiai{
\begin{itemize}
    \item[a)]
    \begin{itemize}
        \item Nam là một trong các thành viên của Chuyên đề 1, vậy Nam là một phần tử của tập hợp $A$. Kí hiệu: $\text{Nam} \in A$.
        \item Ngân là một thành viên của Chuyên đề 1 nhưng không có tên trong danh sách thành viên của Chuyên đề 2. Vậy Ngân không phải là phần tử của tập hợp $B$. Kí hiệu: $\text{Ngân} \notin B$.
    \end{itemize}
    \item[b)] Mô tả tập hợp $A$ và $B$ bằng cách liệt kê các phần tử:
    \[
    \begin{aligned}
        A &= \{\text{Nam}; \text{Hương}; \text{Chi}; \text{Tú}; \text{Bình}; \text{Ngân}; \text{Khánh}\} \\
        B &= \{\text{Hương}; \text{Chi}; \text{Tú}; \text{Khánh}; \text{Bình}; \text{Hân}; \text{Hiền}; \text{Lam}\}.
    \end{aligned}
    \]
\end{itemize}
}

\begin{hoatdong}{Hoạt động 2}
Cho tập hợp: $C = \{\text{châu Á}; \text{châu Âu}; \text{châu Đại Dương}; \text{châu Mĩ}; \text{châu Nam Cực}; \text{châu Phi}\}$.
\begin{itemize}
    \item[a)] Hãy chỉ ra tính chất đặc trưng cho các phần tử của tập hợp $C$.
    \item[b)] Tập hợp $C$ có bao nhiêu phần tử?
\end{itemize}
\end{hoatdong}
\loigiai{
\begin{itemize}
    \item[a)] Tính chất đặc trưng cho các phần tử của tập hợp $C$ là: "là một châu lục trên Trái Đất".
    Vậy, tập hợp $C$ có thể được viết dưới dạng chỉ ra tính chất đặc trưng là:
    $C = \{x \mid x \text{ là một châu lục trên Trái Đất}\}$.
    \item[b)] Bằng cách đếm các phần tử đã liệt kê, ta thấy tập hợp $C$ có $6$ phần tử. Kí hiệu: $n(C) = 6$.
\end{itemize}
}

\begin{kienthuc}
\paragraph{Tập hợp và phần tử}
\begin{itemize}
    \item Một \chuy{tập hợp} là một bộ sưu tập các đối tượng phân biệt. Các đối tượng đó được gọi là các \chuy{phần tử} của tập hợp.
    \item Kí hiệu $a \in S$ để chỉ phần tử $a$ thuộc tập hợp $S$.
    \item Kí hiệu $a \notin S$ để chỉ phần tử $a$ không thuộc tập hợp $S$.
\end{itemize}

\paragraph{Cách mô tả tập hợp}
Có hai cách thường dùng để mô tả một tập hợp:
\begin{itemize}
    \item \chuy{Cách 1: Liệt kê các phần tử của tập hợp.} Các phần tử được viết trong dấu ngoặc nhọn $\{\dots\}$, cách nhau bởi dấu chấm phẩy hoặc dấu phẩy. Mỗi phần tử chỉ được liệt kê một lần và thứ tự liệt kê không quan trọng.
    \item \chuy{Cách 2: Chỉ ra tính chất đặc trưng cho các phần tử của tập hợp.} Tập hợp được viết dưới dạng $S = \{x \mid x \text{ có tính chất } P\}$.
\end{itemize}

\paragraph{Tập rỗng và số phần tử của tập hợp}
\begin{itemize}
    \item Tập hợp không chứa phần tử nào được gọi là \chuy{tập rỗng}, kí hiệu là $\emptyset$.
    \item Số phần tử của tập hợp $S$ được kí hiệu là $n(S)$.
\end{itemize}
\end{kienthuc}

\begin{vd}
Cho tập hợp $D = \{n \in \mathbb{N} \mid n \text{ là số nguyên tố}, 5 < n < 20\}$.
\begin{itemize}
    \item[a)] Trong các số $5; 12; 17; 18$, số nào thuộc tập $D$, số nào không thuộc tập $D$?
    \item[b)] Viết tập hợp $D$ bằng cách liệt kê các phần tử.
\end{itemize}
\end{vd}
\loigiai{
\begin{itemize}
    \item[a)] Để xác định một số có thuộc tập $D$ hay không, ta kiểm tra xem số đó có thỏa mãn cả ba điều kiện: là số tự nhiên, là số nguyên tố, và lớn hơn $5$ nhỏ hơn $20$.
    \begin{itemize}
        \item Số $5$: là số tự nhiên, là số nguyên tố, nhưng $5$ không lớn hơn $5$. Vậy $5 \notin D$.
        \item Số $12$: là số tự nhiên, $5 < 12 < 20$, nhưng $12$ không phải là số nguyên tố (vì $12 = 2^2 \cdot 3$). Vậy $12 \notin D$.
        \item Số $17$: là số tự nhiên, là số nguyên tố, và $5 < 17 < 20$. Vậy $17 \in D$.
        \item Số $18$: là số tự nhiên, $5 < 18 < 20$, nhưng $18$ không phải là số nguyên tố (vì $18 = 2 \cdot 3^2$). Vậy $18 \notin D$.
    \end{itemize}
    \item[b)] Để viết tập hợp $D$ bằng cách liệt kê các phần tử, ta tìm tất cả các số tự nhiên là số nguyên tố và nằm giữa $5$ và $20$.
    Các số nguyên tố lớn hơn $5$ là $7, 11, 13, 17, 19, 23, \dots$.
    Các số nguyên tố nhỏ hơn $20$ là $2, 3, 5, 7, 11, 13, 17, 19$.
    Vậy, các số nguyên tố thỏa mãn $5 < n < 20$ là $7, 11, 13, 17, 19$.
    Do đó, $D = \{7; 11; 13; 17; 19\}$.
\end{itemize}
}

\begin{lt}
Gọi $X$ là tập nghiệm của phương trình $x^2 - 24x + 143 = 0$. Các mệnh đề sau đúng hay sai?
\begin{itemize}
    \item[a)] $13 \in X$;
    \item[b)] $11 \notin X$;
    \item[c)] $n(X) = 2$.
\end{itemize}
\end{lt}
\loigiai{
Trước hết, ta cần tìm tập nghiệm $X$ của phương trình $x^2 - 24x + 143 = 0$.
Ta có thể giải phương trình bậc hai bằng cách phân tích thành nhân tử hoặc dùng công thức nghiệm.
$x^2 - 24x + 143 = 0$
$\Leftrightarrow x^2 - 11x - 13x + 143 = 0$
$\Leftrightarrow x(x - 11) - 13(x - 11) = 0$
$\Leftrightarrow (x - 11)(x - 13) = 0$
$\Leftrightarrow \left[ \begin{aligned} x - 11 &= 0 \\ x - 13 &= 0 \end{aligned} \right.$
$\Leftrightarrow \left[ \begin{aligned} x &= 11 \\ x &= 13 \end{aligned} \right.$
Vậy, tập nghiệm của phương trình là $X = \{11; 13\}$.

Bây giờ, ta xét tính đúng sai của các mệnh đề:
\begin{itemize}
    \item[a)] Mệnh đề "$13 \in X$": Vì $13$ là một nghiệm của phương trình, nên $13$ là một phần tử của tập $X$. Vậy mệnh đề này là \chuy{đúng}.
    \item[b)] Mệnh đề "$11 \notin X$": Vì $11$ là một nghiệm của phương trình, nên $11$ là một phần tử của tập $X$. Do đó, việc nói $11 \notin X$ là sai. Vậy mệnh đề này là \chuy{sai}.
    \item[c)] Mệnh đề "$n(X) = 2$": Tập $X$ có hai phần tử là $11$ và $13$. Vậy số phần tử của $X$ là $2$. Mệnh đề này là \chuy{đúng}.
\end{itemize}
}

\paragraph{Tập hợp con}

\begin{hoatdong}{Hoạt động 3}
Gọi $H$ là tập hợp các bạn tham gia Chuyên đề 2 trong tình huống mở đầu có tên bắt đầu bằng chữ $H$. Các phần tử của tập hợp $H$ có là phần tử của tập hợp $B$ trong HĐ1 không?
\end{hoatdong}
\loigiai{
Trong Hoạt động 1, tập hợp $B$ là tập hợp các thành viên tham gia Chuyên đề 2:
$B = \{\text{Hương}; \text{Chi}; \text{Tú}; \text{Khánh}; \text{Bình}; \text{Hân}; \text{Hiền}; \text{Lam}\}$.

Tập hợp $H$ là tập hợp các bạn tham gia Chuyên đề 2 có tên bắt đầu bằng chữ $H$.
Các thành viên trong $B$ có tên bắt đầu bằng chữ $H$ là: Hương, Hân, Hiền.
Vậy, $H = \{\text{Hương}; \text{Hân}; \text{Hiền}\}$.

Quan sát các phần tử của $H$:
\begin{itemize}
    \item Hương $\in B$
    \item Hân $\in B$
    \item Hiền $\in B$
\end{itemize}
Tất cả các phần tử của tập hợp $H$ đều là phần tử của tập hợp $B$.
}

\begin{kienthuc}
\begin{itemize}
    \item Nếu mọi phần tử của tập hợp $T$ đều là phần tử của tập hợp $S$ thì ta nói $T$ là một \chuy{tập hợp con} (hay \chuy{tập con}) của $S$ và viết là $T \subset S$.
    \item Định nghĩa bằng kí hiệu: $T \subset S \Leftrightarrow (\forall x, x \in T \Rightarrow x \in S)$.
    \item Nếu $T \subset S$ và $T \neq S$ thì ta nói $T$ là \chuy{tập con thực sự} của $S$.
    \item Quy ước tập rỗng $\emptyset$ là tập con của mọi tập hợp.
\end{itemize}
\end{kienthuc}

\begin{haicotchay}{Minh họa tập con}{
\begin{tikzpicture}
\draw (0,0) circle (2cm);
\draw (0.5,0.5) circle (1cm);
\node at (0, 2.2) {$S$};
\node at (0.5, 0.5) {$T$};
\end{tikzpicture}
}
\end{haicotchay}

\begin{vd}
Cho tập hợp $S = \{2; 3; 5\}$. Những tập hợp nào sau đây là tập con của $S$?
\[
S_1 = \{3\}, \quad S_2 = \{0; 2\}, \quad S_3 = \{3; 5\}
\]
\end{vd}
\loigiai{
Để xác định một tập hợp có phải là tập con của $S$ hay không, ta kiểm tra xem tất cả các phần tử của tập hợp đó có thuộc $S$ hay không.
\begin{itemize}
    \item Với $S_1 = \{3\}$: Phần tử $3 \in S$. Vậy $S_1 \subset S$.
    \item Với $S_2 = \{0; 2\}$: Phần tử $0 \notin S$. Vậy $S_2$ không phải là tập con của $S$. Kí hiệu: $S_2 \not\subset S$.
    \item Với $S_3 = \{3; 5\}$: Các phần tử $3 \in S$ và $5 \in S$. Vậy $S_3 \subset S$.
\end{itemize}
}

\begin{lt}
Cho tập hợp $A = \{x \in \mathbb{N} \mid x \text{ là ước của } 12\}$.
\begin{itemize}
    \item[a)] Viết tập hợp $A$ bằng cách liệt kê các phần tử.
    \item[b)] Tập hợp $B = \{1; 2; 3; 4\}$ có phải là tập con của $A$ không?
\end{itemize}
\end{lt}
\loigiai{
\begin{itemize}
    \item[a)] Tập hợp $A$ gồm các số tự nhiên là ước của $12$.
    Các ước của $12$ là $1, 2, 3, 4, 6, 12$.
    Vậy, $A = \{1; 2; 3; 4; 6; 12\}$.
    \item[b)] Để kiểm tra $B$ có phải là tập con của $A$ không, ta xem xét tất cả các phần tử của $B$ có thuộc $A$ không.
    Các phần tử của $B$ là $1, 2, 3, 4$.
    Ta thấy: $1 \in A$, $2 \in A$, $3 \in A$, $4 \in A$.
    Vì mọi phần tử của $B$ đều là phần tử của $A$, nên $B$ là tập con của $A$. Kí hiệu: $B \subset A$.
\end{itemize}
}

\paragraph{Hai tập hợp bằng nhau}

\begin{kienthuc}
Hai tập hợp $A$ và $B$ được gọi là \chuy{bằng nhau} nếu $A$ là tập con của $B$ và $B$ là tập con của $A$. Kí hiệu là $A = B$.
Định nghĩa bằng kí hiệu: $A = B \Leftrightarrow (A \subset B \text{ và } B \subset A)$.
\end{kienthuc}

\begin{vd}
Cho tập hợp $A = \{x \in \mathbb{Z} \mid (x-1)(x^2-4)=0\}$ và $B = \{-2; 1; 2\}$. Hai tập hợp này có bằng nhau không?
\end{vd}
\loigiai{
Trước hết, ta cần xác định các phần tử của tập hợp $A$.
Phương trình $(x-1)(x^2-4)=0$ có các nghiệm:
$x-1=0 \Rightarrow x=1$
$x^2-4=0 \Rightarrow x^2=4 \Rightarrow x=\pm 2$
Vậy, tập nghiệm của phương trình là $A = \{-2; 1; 2\}$.
Ta thấy tập hợp $B$ cũng là $B = \{-2; 1; 2\}$.
Vì $A$ và $B$ có cùng các phần tử, nên $A = B$.
}

\subsubsection{Các phép toán trên tập hợp}

\paragraph{Phép giao của hai tập hợp}

\begin{hoatdong}{Hoạt động 4}
Trong Hoạt động 1, hãy chỉ ra các thành viên tham gia cả Chuyên đề 1 và Chuyên đề 2.
\end{hoatdong}
\loigiai{
Tập hợp $A = \{\text{Nam}; \text{Hương}; \text{Chi}; \text{Tú}; \text{Bình}; \text{Ngân}; \text{Khánh}\}$
Tập hợp $B = \{\text{Hương}; \text{Chi}; \text{Tú}; \text{Khánh}; \text{Bình}; \text{Hân}; \text{Hiền}; \text{Lam}\}$

Các thành viên tham gia cả Chuyên đề 1 và Chuyên đề 2 là những thành viên có mặt trong cả hai danh sách $A$ và $B$.
Bằng cách so sánh hai danh sách, ta tìm được các thành viên chung: Hương, Chi, Tú, Bình, Khánh.
Vậy, các thành viên tham gia cả hai chuyên đề là: Hương, Chi, Tú, Bình, Khánh.
}

\begin{kienthuc}
\begin{itemize}
    \item \chuy{Phép giao} của hai tập hợp $A$ và $B$, kí hiệu là $A \cap B$, là tập hợp gồm các phần tử chung của $A$ và $B$.
    \item Định nghĩa bằng kí hiệu: $A \cap B = \{x \mid x \in A \text{ và } x \in B\}$.
\end{itemize}
\end{kienthuc}

\begin{haicotchay}{Minh họa phép giao}{
\begin{tikzpicture}
\fill[blue!20] (0,0) circle (1.5cm);
\fill[red!20] (2,0) circle (1.5cm);
\fill[green!50!black!50] (intersection of 0,0 circle 1.5 and 2,0 circle 1.5);
\draw (0,0) circle (1.5cm);
\draw (2,0) circle (1.5cm);
\node at (-1, 1) {$A$};
\node at (3, 1) {$B$};
\node at (1, 0) {$A \cap B$};
\end{tikzpicture}
}
\end{haicotchay}

\begin{vd}
Cho $A = \{0; 1; 2; 3; 4\}$ và $B = \{2; 4; 6; 8\}$. Tìm $A \cap B$.
\end{vd}
\loigiai{
Để tìm $A \cap B$, ta liệt kê các phần tử có mặt trong cả hai tập hợp $A$ và $B$.
Các phần tử chung của $A$ và $B$ là $2$ và $4$.
Vậy, $A \cap B = \{2; 4\}$.
}

\paragraph{Phép hợp của hai tập hợp}

\begin{hoatdong}{Hoạt động 5}
Trong Hoạt động 1, hãy chỉ ra các thành viên tham gia ít nhất một trong hai Chuyên đề 1 hoặc Chuyên đề 2.
\end{hoatdong}
\loigiai{
Tập hợp $A = \{\text{Nam}; \text{Hương}; \text{Chi}; \text{Tú}; \text{Bình}; \text{Ngân}; \text{Khánh}\}$
Tập hợp $B = \{\text{Hương}; \text{Chi}; \text{Tú}; \text{Khánh}; \text{Bình}; \text{Hân}; \text{Hiền}; \text{Lam}\}$

Các thành viên tham gia ít nhất một trong hai chuyên đề là tất cả các thành viên có mặt trong $A$ hoặc trong $B$ (hoặc cả hai). Để tìm tập hợp này, ta gộp tất cả các phần tử của $A$ và $B$ lại, mỗi phần tử chỉ liệt kê một lần.
Các phần tử của $A$: Nam, Hương, Chi, Tú, Bình, Ngân, Khánh.
Các phần tử của $B$: Hương, Chi, Tú, Khánh, Bình, Hân, Hiền, Lam.
Gộp lại, ta được: Nam, Hương, Chi, Tú, Bình, Ngân, Khánh, Hân, Hiền, Lam.
Vậy, tập hợp các thành viên tham gia ít nhất một trong hai chuyên đề là:
$\{\text{Nam}; \text{Hương}; \text{Chi}; \text{Tú}; \text{Bình}; \text{Ngân}; \text{Khánh}; \text{Hân}; \text{Hiền}; \text{Lam}\}$.
}

\begin{kienthuc}
\begin{itemize}
    \item \chuy{Phép hợp} của hai tập hợp $A$ và $B$, kí hiệu là $A \cup B$, là tập hợp gồm tất cả các phần tử thuộc $A$ hoặc thuộc $B$ (hoặc thuộc cả hai).
    \item Định nghĩa bằng kí hiệu: $A \cup B = \{x \mid x \in A \text{ hoặc } x \in B\}$.
\end{itemize}
\end{kienthuc}

\begin{haicotchay}{Minh họa phép hợp}{
\begin{tikzpicture}
\fill[blue!20] (0,0) circle (1.5cm);
\fill[red!20] (2,0) circle (1.5cm);
\draw (0,0) circle (1.5cm);
\draw (2,0) circle (1.5cm);
\node at (-1, 1) {$A$};
\node at (3, 1) {$B$};
\node at (1, 0) {$A \cup B$};
\end{tikzpicture}
}
\end{haicotchay}

\begin{vd}
Cho $A = \{0; 1; 2; 3; 4\}$ và $B = \{2; 4; 6; 8\}$. Tìm $A \cup B$.
\end{vd}
\loigiai{
Để tìm $A \cup B$, ta gộp tất cả các phần tử của $A$ và $B$ lại, mỗi phần tử chỉ liệt kê một lần.
Các phần tử của $A$: $0, 1, 2, 3, 4$.
Các phần tử của $B$: $2, 4, 6, 8$.
Gộp lại, ta được: $0, 1, 2, 3, 4, 6, 8$.
Vậy, $A \cup B = \{0; 1; 2; 3; 4; 6; 8\}$.
}

\paragraph{Phép hiệu của hai tập hợp}

\begin{hoatdong}{Hoạt động 6}
Trong Hoạt động 1, hãy chỉ ra các thành viên tham gia Chuyên đề 1 nhưng không tham gia Chuyên đề 2.
\end{hoatdong}
\loigiai{
Tập hợp $A = \{\text{Nam}; \text{Hương}; \text{Chi}; \text{Tú}; \text{Bình}; \text{Ngân}; \text{Khánh}\}$
Tập hợp $B = \{\text{Hương}; \text{Chi}; \text{Tú}; \text{Khánh}; \text{Bình}; \text{Hân}; \text{Hiền}; \text{Lam}\}$

Các thành viên tham gia Chuyên đề 1 nhưng không tham gia Chuyên đề 2 là những thành viên thuộc $A$ nhưng không thuộc $B$.
Ta loại bỏ các phần tử chung của $A$ và $B$ ra khỏi $A$.
Các phần tử chung là: Hương, Chi, Tú, Bình, Khánh.
Các phần tử của $A$ không thuộc $B$ là: Nam, Ngân.
Vậy, các thành viên tham gia Chuyên đề 1 nhưng không tham gia Chuyên đề 2 là: Nam, Ngân.
}

\begin{kienthuc}
\begin{itemize}
    \item \chuy{Phép hiệu} của hai tập hợp $A$ và $B$, kí hiệu là $A \setminus B$ (hoặc $A - B$), là tập hợp gồm các phần tử thuộc $A$ nhưng không thuộc $B$.
    \item Định nghĩa bằng kí hiệu: $A \setminus B = \{x \mid x \in A \text{ và } x \notin B\}$.
\end{itemize}
\end{kienthuc}

\begin{haicotchay}{Minh họa phép hiệu}{
\begin{tikzpicture}
\fill[blue!20] (0,0) circle (1.5cm);
\fill[white] (2,0) circle (1.5cm);
\fill[white] (intersection of 0,0 circle 1.5 and 2,0 circle 1.5);
\draw (0,0) circle (1.5cm);
\draw (2,0) circle (1.5cm);
\node at (-0.7, 0.7) {$A$};
\node at (2.2, 0.7) {$B$};
\node at (-0.7, 0) {$A \setminus B$};
\end{tikzpicture}
}
\end{haicotchay}

\begin{vd}
Cho $A = \{0; 1; 2; 3; 4\}$ và $B = \{2; 4; 6; 8\}$. Tìm $A \setminus B$ và $B \setminus A$.
\end{vd}
\loigiai{
\begin{itemize}
    \item Để tìm $A \setminus B$, ta tìm các phần tử thuộc $A$ nhưng không thuộc $B$.
    Các phần tử của $A$ là $0, 1, 2, 3, 4$.
    Các phần tử chung của $A$ và $B$ là $2, 4$.
    Loại bỏ $2, 4$ khỏi $A$, ta còn lại $0, 1, 3$.
    Vậy $A \setminus B = \{0; 1; 3\}$.
    \item Để tìm $B \setminus A$, ta tìm các phần tử thuộc $B$ nhưng không thuộc $A$.
    Các phần tử của $B$ là $2, 4, 6, 8$.
    Các phần tử chung của $A$ và $B$ là $2, 4$.
    Loại bỏ $2, 4$ khỏi $B$, ta còn lại $6, 8$.
    Vậy $B \setminus A = \{6; 8\}$.
\end{itemize}
}

\paragraph{Phần bù của một tập hợp}

\begin{kienthuc}
\begin{itemize}
    \item Cho $A$ là một tập con của tập hợp $E$. \chuy{Phần bù} của $A$ trong $E$, kí hiệu là $C_E A$ (hoặc $A'$ nếu tập $E$ được hiểu ngầm), là tập hợp gồm các phần tử thuộc $E$ nhưng không thuộc $A$.
    \item Định nghĩa bằng kí hiệu: $C_E A = \{x \mid x \in E \text{ và } x \notin A\}$.
    \item Từ định nghĩa, ta thấy $C_E A = E \setminus A$.
\end{itemize}
\end{kienthuc}

\begin{haicotchay}{Minh họa phần bù}{
\begin{tikzpicture}
\draw (0,0) rectangle (4,3);
\fill[blue!20] (0,0) rectangle (4,3);
\fill[white] (2,1.5) circle (1);
\node at (3.5, 2.5) {$E$};
\node at (2, 1.5) {$A$};
\node at (0.5, 0.5) {$C_E A$};
\end{tikzpicture}
}
\end{haicotchay}

\begin{vd}
Cho $E = \{0; 1; 2; 3; 4; 5; 6; 7; 8; 9\}$ và $A = \{0; 2; 4; 6; 8\}$. Tìm $C_E A$.
\end{vd}
\loigiai{
Để tìm $C_E A$, ta tìm các phần tử thuộc $E$ nhưng không thuộc $A$.
Các phần tử của $E$ là $0, 1, 2, 3, 4, 5, 6, 7, 8, 9$.
Các phần tử của $A$ là $0, 2, 4, 6, 8$.
Loại bỏ các phần tử của $A$ ra khỏi $E$, ta còn lại $1, 3, 5, 7, 9$.
Vậy $C_E A = \{1; 3; 5; 7; 9\}$.
}

\subsubsection{Các tập hợp số}

\begin{kienthuc}
\paragraph{Các tập hợp số cơ bản}
\begin{itemize}
    \item Tập hợp số tự nhiên: $\mathbb{N} = \{0; 1; 2; 3; \dots\}$
    \item Tập hợp số tự nhiên khác 0: $\mathbb{N}^* = \{1; 2; 3; \dots\}$
    \item Tập hợp số nguyên: $\mathbb{Z} = \{\dots; -2; -1; 0; 1; 2; \dots\}$
    \item Tập hợp số hữu tỉ: $\mathbb{Q} = \left\{\dfrac{a}{b} \mid a \in \mathbb{Z}, b \in \mathbb{Z}, b \neq 0\right\}$
    \item Tập hợp số thực: $\mathbb{R}$
\end{itemize}
Mối quan hệ bao hàm giữa các tập hợp số này là: $\mathbb{N} \subset \mathbb{Z} \subset \mathbb{Q} \subset \mathbb{R}$.
\end{kienthuc}

\paragraph{Các tập con thường dùng của tập số thực}

\begin{kienthuc}
Cho $a, b$ là hai số thực với $a < b$.
\begin{itemize}
    \item \chuy{Khoảng}:
    \begin{itemize}
        \item $(a; b) = \{x \in \mathbb{R} \mid a < x < b\}$
        \item $(a; +\infty) = \{x \in \mathbb{R} \mid x > a\}$
        \item $(-\infty; b) = \{x \in \mathbb{R} \mid x < b\}$
        \item $(-\infty; +\infty) = \mathbb{R}$
    \end{itemize}
    \item \chuy{Đoạn}:
    \begin{itemize}
        \item $[a; b] = \{x \in \mathbb{R} \mid a \le x \le b\}$
    \end{itemize}
    \item \chuy{Nửa khoảng}:
    \begin{itemize}
        \item $[a; b) = \{x \in \mathbb{R} \mid a \le x < b\}$
        \item $(a; b] = \{x \in \mathbb{R} \mid a < x \le b\}$
        \item $[a; +\infty) = \{x \in \mathbb{R} \mid x \ge a\}$
        \item $(-\infty; b] = \{x \in \mathbb{R} \mid x \le b\}$
    \end{itemize}
\end{itemize}
\end{kienthuc}

\begin{vd}
Hãy xác định các tập hợp sau và biểu diễn chúng trên trục số:
\begin{itemize}
    \item[a)] $(-2; 5] \cap [1; 7)$;
    \item[b)] $(-\infty; 3) \cup [0; 4]$;
    \item[c)] $[-1; 6) \setminus (2; 8)$.
\end{itemize}
\end{vd}
\loigiai{
\begin{enumerate}
    \item[a)] Xác định $(-2; 5] \cap [1; 7)$:
    Tập hợp $(-2; 5]$ bao gồm các số thực $x$ sao cho $-2 < x \le 5$.
    Tập hợp $[1; 7)$ bao gồm các số thực $x$ sao cho $1 \le x < 7$.
    Phép giao là tập hợp các phần tử chung, tức là các số $x$ thỏa mãn cả hai điều kiện:
    $1 \le x \le 5$.
    Vậy, $(-2; 5] \cap [1; 7) = [1; 5]$.
    \begin{center}
    \begin{tikzpicture}
    \draw[->] (-3,0) -- (8,0) node[right] {$x$};
    \draw (-2,0) node[below] {$-2$} circle (2pt);
    \draw (5,0) node[below] {$5$} circle (2pt);
    \draw[line width=1.5pt, blue, opacity=0.6] (-2,0) -- (5,0);
    \draw (1,0) node[below] {$1$} circle (2pt);
    \draw (7,0) node[below] {$7$} circle (2pt);
    \draw[line width=1.5pt, red, opacity=0.6] (1,0) -- (7,0);
    \draw[line width=2.5pt, green!50!black] (1,0) -- (5,0);
    \fill (1,0) circle (2pt);
    \fill (5,0) circle (2pt);
    \node[above] at (1,0) {$[$};
    \node[above] at (5,0) {$]$};
    \node[above] at (-2,0) {$( $};
    \node[above] at (7,0) {$)$};
    \end{tikzpicture}
    \end{center}
    \item[b)] Xác định $(-\infty; 3) \cup [0; 4]$:
    Tập hợp $(-\infty; 3)$ bao gồm các số thực $x$ sao cho $x < 3$.
    Tập hợp $[0; 4]$ bao gồm các số thực $x$ sao cho $0 \le x \le 4$.
    Phép hợp là tập hợp các phần tử thuộc ít nhất một trong hai tập hợp.
    Khi hợp hai tập này, ta lấy từ giá trị nhỏ nhất đến giá trị lớn nhất của các khoảng.
    Các số $x$ thỏa mãn $x < 3$ hoặc $0 \le x \le 4$.
    Điều này tương đương với $x \le 4$.
    Vậy, $(-\infty; 3) \cup [0; 4] = (-\infty; 4]$.
    \begin{center}
    \begin{tikzpicture}
    \draw[->] (-3,0) -- (5,0) node[right] {$x$};
    \draw (3,0) node[below] {$3$} circle (2pt);
    \draw[line width=1.5pt, blue, opacity=0.6] (-3,0) -- (3,0);
    \draw (0,0) node[below] {$0$} circle (2pt);
    \draw (4,0) node[below] {$4$} circle (2pt);
    \draw[line width=1.5pt, red, opacity=0.6] (0,0) -- (4,0);
    \draw[line width=2.5pt, green!50!black] (-3,0) -- (4,0);
    \fill (4,0) circle (2pt);
    \node[above] at (3,0) {$)$};
    \node[above] at (0,0) {$[$};
    \node[above] at (4,0) {$]$};
    \end{tikzpicture}
    \end{center}
    \item[c)] Xác định $[-1; 6) \setminus (2; 8)$:
    Tập hợp $[-1; 6)$ bao gồm các số thực $x$ sao cho $-1 \le x < 6$.
    Tập hợp $(2; 8)$ bao gồm các số thực $x$ sao cho $2 < x < 8$.
    Phép hiệu là tập hợp các phần tử thuộc $[-1; 6)$ nhưng không thuộc $(2; 8)$.
    Tức là, ta lấy các số $x$ thỏa mãn $-1 \le x < 6$ và $x \notin (2; 8)$.
    Điều kiện $x \notin (2; 8)$ có nghĩa là $x \le 2$ hoặc $x \ge 8$.
    Kết hợp với $-1 \le x < 6$:
    \begin{itemize}
        \item Nếu $x \le 2$: ta có $-1 \le x \le 2$.
        \item Nếu $x \ge 8$: không có giá trị nào thỏa mãn đồng thời $x < 6$ và $x \ge 8$.
    \end{itemize}
    Vậy, $[-1; 6) \setminus (2; 8) = [-1; 2]$.
    \begin{center}
    \begin{tikzpicture}
    \draw[->] (-2,0) -- (9,0) node[right] {$x$};
    \draw (-1,0) node[below] {$-1$} circle (2pt);
    \draw (6,0) node[below] {$6$} circle (2pt);
    \draw[line width=1.5pt, blue, opacity=0.6] (-1,0) -- (6,0);
    \draw (2,0) node[below] {$2$} circle (2pt);
    \draw (8,0) node[below] {$8$} circle (2pt);
    \draw[line width=1.5pt, red, opacity=0.6] (2,0) -- (8,0);
    \draw[line width=2.5pt, green!50!black] (-1,0) -- (2,0);
    \fill (-1,0) circle (2pt);
    \fill (2,0) circle (2pt);
    \node[above] at (-1,0) {$[$};
    \node[above] at (6,0) {$)$};
    \node[above] at (2,0) {$($};
    \node[above] at (8,0) {$)$};
    \end{tikzpicture}
    \end{center}
\end{enumerate}
}

\begin{emdahoc}
\begin{itemize}
    \item Khái niệm tập hợp, phần tử thuộc tập hợp, tập rỗng, tập hợp con, hai tập hợp bằng nhau và các cách mô tả tập hợp (liệt kê phần tử, chỉ ra tính chất đặc trưng).
    \item Định nghĩa và kí hiệu về các phép toán giao, hợp, hiệu của hai tập hợp và phần bù của một tập con trong tập lớn hơn.
    \item Mối quan hệ giữa các tập hợp số tự nhiên ($\mathbb{N}$), số nguyên ($\mathbb{Z}$), số hữu tỉ ($\mathbb{Q}$), số thực ($\mathbb{R}$) và các tập con thường dùng của tập số thực (khoảng, đoạn, nửa khoảng).
\end{itemize}
\end{emdahoc}

\begin{emcothe}
\begin{itemize}
    \item Biểu diễn linh hoạt các tập hợp bằng cách liệt kê phần tử hoặc chỉ ra tính chất đặc trưng.
    \item Thực hiện thành thạo các phép toán giao, hợp, hiệu trên các khoảng, đoạn, nửa khoảng của tập số thực và biểu diễn chúng chính xác trên trục số.
    \item Giải quyết tốt các bài toán đếm thực tế (sử dụng biểu đồ Venn) liên quan đến các nhóm đối tượng giao thoa nhau trong đời sống xã hội.
\end{itemize}
\end{emcothe}

\begin{emcobiet}
\paragraph{John Venn và Biểu đồ Venn}
Khái niệm tập hợp và các phép toán trên tập hợp được mô tả trực quan sinh động nhất thông qua một công cụ đồ họa hình học vô cùng nổi tiếng gọi là \chuy{Biểu đồ Venn} (Venn diagram). Công cụ này được giới thiệu lần đầu tiên vào năm 1880 bởi nhà logic học và triết học người Anh \chuy{John Venn} (1834 – 1923) trong bài báo có tiêu đề "Về sự biểu diễn cơ học và biểu đồ của các mệnh đề và lý luận". Biểu đồ Venn sử dụng các đường cong khép kín (thường là hình tròn hoặc hình elip) nằm trong một hình chữ nhật lớn đại diện cho vũ trụ tập hợp để mô tả các mối quan hệ logic giữa các tập hợp như giao, hợp, hiệu và phần bù. Không chỉ dừng lại ở toán học thuần túy, biểu đồ Venn ngày nay đã trở thành một công cụ tư duy không thể thiếu trong giảng dạy logic, xác suất thống kê, khoa học máy tính và thậm chí cả trong các chiến lược tiếp thị, phân tích dữ liệu kinh doanh của các doanh nghiệp hiện đại.
\end{emcobiet}
